\documentclass[a4paper,fleqn]{cas-dc}

\usepackage[numbers,sort&compress]{natbib}
\usepackage{graphicx}

\begin{document}

\title[mode=title]{Self-Supervised Pretraining for Defect Detection on Production Lines}

\author[1]{Wei Chen}
\cormark[1]
\ead{wei.chen@nthu.edu.tw}

\author[1]{Ming-Hui Lee}

\author[2]{Anna Schmidt}

\affiliation[1]{organization={National Tsing Hua University},
  addressline={101 Sec. 2 Kuang-Fu Rd.},
  city={Hsinchu},
  postcode={30013},
  country={Taiwan}}

\affiliation[2]{organization={Technical University of Munich},
  addressline={Arcisstrasse 21},
  city={Munich},
  postcode={80333},
  country={Germany}}

\cortext[1]{Corresponding author}

\begin{abstract}
Defect detection on production lines is a long-standing challenge where
labeled data is scarce and defect distributions shift over time. We propose
a self-supervised pretraining strategy tailored to industrial imagery,
which learns transferable representations from unlabeled production-line
footage and transfers them to downstream defect classification with minimal
fine-tuning. Our approach combines masked image modeling with a contrastive
objective that respects the temporal structure of conveyor-belt sequences.
We evaluate on four industrial defect datasets spanning metal casting,
textile weaving, printed circuit board inspection, and pharmaceutical
packaging. Our pretrained model consistently outperforms ImageNet-pretrained
baselines by 3 to 6 points in F1 score under low-label regimes, and matches
fully supervised performance using only 10 percent of the labeled data.
We further demonstrate that the learned representations transfer across
defect categories with minimal performance loss, suggesting that
production-line pretraining captures factory-specific visual regularities
that generic pretraining does not. Ablations isolate the contribution of
the temporal contrastive objective and show that it is the key ingredient
enabling cross-domain transfer within the industrial setting.
\end{abstract}

\begin{keywords}
self-supervised learning \sep defect detection \sep industrial vision \sep transfer learning
\end{keywords}

\begin{highlights}
\item Self-supervised pretraining tailored to industrial production-line imagery
\item Combines masked image modeling with temporal contrastive objective
\item Evaluated on four industrial defect datasets spanning diverse domains
\item Matches supervised performance with only 10 percent labeled data
\item Temporal contrastive objective is key to cross-domain transfer
\end{highlights}

\section{Introduction}

Defect detection on production lines is a core smart-manufacturing task
\citep{smith2020,doe2021}. See Figure~\ref{fig:pipeline} for our approach.

\section{Related Work}

Self-supervised methods have transformed computer vision \citep{jones2022}.
Industrial applications remain comparatively underexplored \citep{brown2023}.

\section{Method}

Figure~\ref{fig:pipeline} shows the pretraining pipeline. The contrastive
loss is
\begin{equation}
L_\text{con} = -\log \frac{\exp(s(z_i, z_j) / \tau)}{\sum_k \exp(s(z_i, z_k) / \tau)}.
\label{eq:contrastive}
\end{equation}

\begin{figure}
\includegraphics[width=\linewidth]{pipeline.pdf}
\caption{Pretraining pipeline combining masked image modeling with temporal contrastive learning.}
\label{fig:pipeline}
\end{figure}

\section{Experiments}

Quantitative results are summarized in Figure~\ref{fig:results} and
Table~\ref{tab:main}.

\begin{figure}
\includegraphics[width=\linewidth]{results.pdf}
\caption{Downstream F1 scores under varying label budgets.}
\label{fig:results}
\end{figure}

\begin{table}
\caption{Main results across four industrial datasets.}
\label{tab:main}
\begin{tabular}{lcccc}
\hline
Method & Casting & Textile & PCB & Packaging \\
\hline
ImageNet & 0.72 & 0.68 & 0.75 & 0.70 \\
Ours & 0.78 & 0.74 & 0.81 & 0.76 \\
\hline
\end{tabular}
\end{table}

\section{Conclusion}

We have shown that self-supervised pretraining on production-line imagery
consistently improves downstream defect detection. The proposed temporal
contrastive objective is the key ingredient for cross-domain transfer.

\section*{Acknowledgements}

We thank our industrial partners for providing the production-line footage
and the reviewers for their constructive feedback.

\section*{CRediT authorship contribution statement}

\textbf{Wei Chen:} Conceptualization, Methodology, Software, Validation, Writing -- original draft.
\textbf{Ming-Hui Lee:} Data curation, Investigation, Visualization.
\textbf{Anna Schmidt:} Supervision, Writing -- review \& editing, Funding acquisition.

\section*{Declaration of competing interest}

The authors declare that they have no known competing financial interests
or personal relationships that could have appeared to influence the work
reported in this paper.

\section*{Data availability}

The MVTec AD dataset is publicly available. The in-house production-line
footage cannot be shared due to confidentiality agreements with our
industrial partners.

\section*{Declaration of generative AI use}

During the preparation of this work the authors used ChatGPT in order to
polish the language of early drafts. After using this tool, the authors
reviewed and edited the content as needed and take full responsibility
for the content of the publication.

\bibliographystyle{cas-model2-names}
\bibliography{refs}

\end{document}
